\documentclass[11pt]{article}
\usepackage[margin=1in]{geometry}
\usepackage{amsmath,amssymb,amsthm,mathtools}
\usepackage{graphicx}
\usepackage[hidelinks]{hyperref}
\emergencystretch=3em

\newtheorem{theorem}{Theorem}
\newtheorem{lemma}{Lemma}
\newtheorem{remark}{Remark}

\newcommand{\R}{\mathbb R}
\newcommand{\tr}{\operatorname{tr}}
\newcommand{\dd}{\,\mathrm d}

\title{A Trace-Free Slice of the Minor-Invariant Projection Question}
\author{Run fa\_banach\_001, lane 19}
\date{July 19, 2026}

\begin{document}
\maketitle

\section*{Source Question}

Guerra and Raita ask in Question 5.17 of arXiv:1909.03923 for conditions on a
projection
\[
  T:\R^{m\times n}\to V\subseteq \R^{m\times n}
\]
under which one can find a nonzero linear combination $H_s$ of $s$-minors,
$2\le s<\min\{m,n\}$, satisfying
\[
  H_s(X)=H_s(X+Y)
  \quad\text{for every }X\text{ and every }Y\text{ with }T(Y)=0.
\]
Equivalently, $H_s$ is a linear combination of $s$-minors depending only on the
coordinates of $V$. The source passage is shown in Figure~\ref{fig:source}.

\begin{figure}[t]
\centering
\includegraphics[height=.77\textheight,keepaspectratio]{figures/open_question_page.png}
\caption{Question 5.17 in the source paper, printed page 24.}
\label{fig:source}
\end{figure}

This packet solves the natural trace-free square-matrix slice of the question.

\section*{Definitions}

For subsets $I,J\subset\{1,\dots,n\}$ with $|I|=|J|=s$, write
\[
  M_{I,J}(X)=\det X_{I,J},
\]
where $X_{I,J}$ is the submatrix with rows $I$ and columns $J$, both listed in
increasing order. Every linear combination of $s$-minors on $\R^{n\times n}$
has the form
\[
  H(X)=\sum_{|I|=|J|=s}c_{I,J}M_{I,J}(X).
\]

Let
\[
  T(X)=X-\frac{\tr X}{n}I_n.
\]
Then $T$ is the orthogonal projection onto the trace-free matrices and
$\ker T=\R I_n$. Thus the source condition becomes
\[
  H(X+tI_n)=H(X)\quad\text{for all }X\in\R^{n\times n},\ t\in\R.
\]

\section*{Proof Intuition}

If $2s\le n$, there are disjoint row and column sets $I,J$ of size $s$; the
minor $M_{I,J}$ does not involve any diagonal entry, so adding $tI_n$ leaves it
unchanged.

The converse is the useful part. If $2s>n$, every two $s$-subsets $I,J$ meet.
Differentiating the identity $H(X+tI_n)=H(X)$ in $t$ gives a linear relation
among $(s-1)$-minors. The coefficient of a fixed $(s-1)$-minor is a signed sum
of the coefficients of the $s$-minors obtained by adding the same index to the
row and column sets. A descending induction on $|I\cup J|$ then isolates and
kills all coefficients.

\section*{Main Partial Theorem}

\begin{theorem}[Trace-free scalar-shift threshold]
\label{thm:main}
Let $n\ge3$ and $2\le s<n$. Let $T:\R^{n\times n}\to\{X:\tr X=0\}$ be
\[
  T(X)=X-\frac{\tr X}{n}I_n.
\]
There is a nonzero linear combination $H$ of $s$-minors such that
\[
  H(X)=H(X+Y)\quad\text{whenever }T(Y)=0
\]
if and only if
\[
  2s\le n.
\]
\end{theorem}

\begin{proof}
Since $\ker T=\R I_n$, the invariance condition is exactly
\[
  H(X+tI_n)=H(X)\quad\text{for all }X,t.
\]

If $2s\le n$, choose disjoint subsets $I,J\subset\{1,\dots,n\}$ with
$|I|=|J|=s$. For $i\in I$ and $j\in J$ we have $i\ne j$, so the entries of the
submatrix $X_{I,J}$ are unchanged by adding $tI_n$. Hence
$H=M_{I,J}$ is a nonzero invariant linear combination of $s$-minors.

Conversely, assume $2s>n$ and let
\[
  H(X)=\sum_{|I|=|J|=s}c_{I,J}M_{I,J}(X)
\]
be invariant. Differentiating at $t=0$ gives
\[
  0=\left.\frac{\dd}{\dd t}\right|_{t=0}H(X+tI_n)
\]
as a polynomial identity in $X$. For a single minor,
\[
  \left.\frac{\dd}{\dd t}\right|_{t=0}M_{I,J}(X+tI_n)
  =
  \sum_{a\in I\cap J}\varepsilon(a;I,J)M_{I\setminus\{a\},J\setminus\{a\}}(X),
\]
where each $\varepsilon(a;I,J)$ is $1$ or $-1$. The precise sign will not
matter.

The minors of a fixed order are linearly independent as polynomials, since a
monomial occurring in $M_{A,B}$ uses exactly the row set $A$ and the column set
$B$. Therefore, for every pair $A,B$ of $(s-1)$-subsets, the coefficient of
$M_{A,B}$ in the differentiated identity is zero:
\[
  \sum_{a\notin A\cup B}
  \varepsilon(a;A\cup\{a\},B\cup\{a\})
  c_{A\cup\{a\},B\cup\{a\}}=0.
  \tag{1}
\]

We prove that all coefficients $c_{I,J}$ vanish by descending induction on
$N=|I\cup J|$. Since $2s>n$, every pair of $s$-sets $I,J$ has
$I\cap J\ne\emptyset$.

Start with $N=n$. Pick $a\in I\cap J$ and set
$A=I\setminus\{a\}$, $B=J\setminus\{a\}$. Then
$A\cup B=(I\cup J)\setminus\{a\}$ has complement exactly $\{a\}$. Thus (1)
contains only the coefficient $c_{I,J}$, up to a nonzero sign, and so
$c_{I,J}=0$.

Now suppose all coefficients with union size larger than $N$ vanish, and take
$I,J$ with $|I\cup J|=N$. Choose $a\in I\cap J$ and set
$A=I\setminus\{a\}$, $B=J\setminus\{a\}$. In (1), the term with the added index
$a$ is precisely $c_{I,J}$, up to sign. Any other possible added index
$b\notin A\cup B$ lies outside $I\cup J$, and the corresponding coefficient
$c_{A\cup\{b\},B\cup\{b\}}$ has union size $N+1$; it is zero by the induction
hypothesis. Hence $c_{I,J}=0$.

The descending induction reaches every pair $I,J$, so all coefficients vanish.
Thus no nonzero invariant combination exists when $2s>n$.
\end{proof}

\section*{Scope and Novelty Check}

The theorem gives an exact existence threshold for the projection onto
trace-free square matrices. It is a partial answer to Question 5.17, not a
classification for arbitrary projections $T$. In the positive range
$2s\le n$, the theorem gives explicit invariants but does not describe the
entire invariant subspace.

Cheap run indexes were searched for arXiv:1909.03923, trace-free, traceless,
scalar-shift, identity direction, minors, and the source question keywords. No
duplicate packet was found. A bounded web search on July 19, 2026 found the
source paper and preprint records but no exact answer to this trace-free slice.

\section*{Computational Sanity Check}

The script \texttt{code/check\_tracefree\_threshold.py} builds the exact
rational matrix of the identity-direction derivative on $s$-minors and checks
the kernel dimensions for $3\le n\le7$. It confirms the threshold above for the
tested cases. This computation is only a verification aid; the proof above is
independent.

\section*{References}

\begin{enumerate}
\item A. Guerra and B. Raita,
\emph{Quasiconvexity, null Lagrangians, and Hardy space integrability under
constant rank constraints}, arXiv:1909.03923; Archive for Rational Mechanics
and Analysis 245 (2022), 279--320.
\end{enumerate}

\end{document}
